\subsection{Note on a Corrected Evaluation Pipeline}

An earlier version of this evaluation pipeline contained two bugs,
corrected prior to the results reported in Section~\ref{sec:results}.
First, the closed/open answer-type label was read from a field that does
not exist on the Hugging Face mirror of VQA-RAD used here, causing every
example to be silently mislabeled open-ended and closed-question accuracy
to be permanently reported as 0.0000. Second, the evaluation set was
constructed by internally re-splitting VQA-RAD's official training
portion (yielding a 180-example test set) rather than using VQA-RAD's
own official 451-example held-out test split, which is what other
published VQA-RAD baselines evaluate against. Both are fixed in the
pipeline this paper's numbers are drawn from: answer type is recovered
from the answer text (Section~\ref{sec:datasets}), and the official test
split is used directly. We document this correction explicitly, rather
than silently revising the numbers, in the interest of transparent
research practice.

\subsection{Dataset Scope}

Both datasets used here are deliberately non-credentialed
(Section~\ref{sec:datasets}), which was necessary for a timeline-constrained
project but has real consequences for generalizability. VQA-RAD's test
split is small (on the order of a few hundred examples), which is why
Section~\ref{sec:experimental_setup} treats statistical significance as a
requirement rather than an optional check, but a small test split also
limits the granularity of any subgroup analysis (e.g.\ by modality or
anatomical region) that a larger benchmark would support. ROCOv2's captions
are figure captions drawn from published literature, not full clinical
reports, and are English-only; neither dataset was designed as a matched
pair, so the retrieval corpus and evaluation benchmark may differ somewhat
in imaging style, patient population, or reporting convention.
MIMIC-CXR, the credentialed dataset most published medical VLM work uses,
is a natural extension once credentialed access is obtained
(Section~\ref{sec:conclusion}).

\subsection{Retriever Design}

The retriever embeds and matches only the query \emph{image}: the CLIP
image encoder retrieves corpus entries by visual similarity alone, and the
VQA question text never participates in retrieval (Section~\ref{sec:method}).
As discussed in Section~\ref{sec:discussion}, this is a plausible
alternative explanation (alongside evidence dilution) for why accuracy
declines as $k$ increases -- a question-aware or cross-modal retriever,
which this work does not implement, could retrieve evidence relevant to
what is specifically being asked rather than only to what the image
visually resembles. The retriever also uses a single, off-the-shelf CLIP
checkpoint (\texttt{openai/clip-vit-base-patch32}) trained on
general-domain web images, not fine-tuned on medical imagery specifically.
Whether a medical-domain-adapted image encoder would retrieve more
clinically relevant evidence, and thereby change the retrieval-augmented
result, is untested here. Likewise, only an exact flat inner-product index
is used; this is an appropriate, deliberate choice at ROCOv2's corpus scale
($\sim$90K vectors, Section~\ref{sec:method}) but would not scale without
revisiting approximate indexing (IVF/HNSW) at a larger corpus size, and the
top-$k$ sweep in Section~\ref{sec:results} covers only a handful of
discrete values, not a continuous sensitivity analysis.

\subsection{Evaluation Methodology}

Closed-question accuracy in this paper relies on an answer-type label
recovered from the answer text (Section~\ref{sec:datasets}: a normalized
answer of exactly ``yes'' or ``no'' is treated as closed-ended), since
the VQA-RAD mirror used here does not expose the original annotation.
This recovers the large majority of the true closed/open split but will
mislabel a small number of closed-ended, non-yes/no questions as
open-ended; the reported closed-question accuracy should be read with
this approximation in mind. Open-ended question scoring relies on
BLEU-4 and ROUGE-L, both surface n-gram overlap metrics that can penalize
a correct answer phrased differently from the reference and, conversely,
reward superficial lexical overlap without semantic correctness. BLEU-4
in particular is close to structurally uninformative at VQA-RAD's typical
open-ended answer length (often 3--4 tokens): even two identical strings
can score well below 1.0 under BLEU-4 with smoothing, since a short
answer offers few or no 4-gram matches regardless of correctness; the
BLEU-4 deltas reported in Section~\ref{sec:results} should be weighted
accordingly relative to ROUGE-L and exact match. The 0.5 ROUGE-L
threshold used to bucket open-ended examples into
\emph{improved}/\emph{regressed} categories for qualitative review
(Section~\ref{sec:experimental_setup}) is a convenience heuristic for
categorization, not a validated correctness threshold, and is not used
anywhere in the headline reported metrics. No clinician review of model
answers for clinical adequacy (as opposed to lexical match against a
reference) was performed; this is a meaningful gap between what this
evaluation measures and what would be required before any claim of
clinical utility.

\subsection{Compute Constraints}

As a resource-constrained, timeline-bound project, this work evaluates a
single baseline checkpoint family
(\texttt{Salesforce/blip2-opt-2.7b}, with
\texttt{microsoft/llava-med-v1.5-mistral-7b} as an alternate candidate) at
a single scale, rather than an ablation across multiple model sizes or
architectures. Similarly, only one prompt template
(\texttt{configs/model\_config.yaml}'s \texttt{rag.evidence\_instruction})
was used for the retrieval-augmented condition; prompt phrasing is known in
the broader RAG literature to affect how well a model makes use of
retrieved context, and this was not itself an experimental variable here.

\subsection{Scope of Claims}

This is a research prototype intended to answer a narrow empirical
question --- does retrieval augmentation change VLM accuracy on a specific
medical VQA benchmark, under one model and retriever configuration --- and
is explicitly not a validated diagnostic or clinical decision-support tool.
Any result reported in Section~\ref{sec:results} should be read as evidence
about this specific pipeline and benchmark, not as a general claim about
retrieval-augmented medical VLMs as a class.
